% 05-experiences-methode (fr): le protocole des expériences, tel que le code l'exécute.
\chapter{Les expériences : protocole et méthodes}\label{chap:05-experiences-methode}

Le sondage des \cref{chap:03-sondage-methode,chap:04-sondage-resultats} dit ce que les articles publient. Les
expériences mettent les failles de 2018 à l'épreuve sur des entraînements réels, faits pour ce rapport. Elles
reprennent le terrain de \textcite{henderson2018matters} : des corps simulés dans le simulateur physique MuJoCo,
que des algorithmes d'apprentissage par renforcement apprennent à faire avancer. Elles le font avec les
algorithmes de 2026 et avec les règles que Joelle Pineau demandait : assez de graines aléatoires, des
intervalles de confiance, des références traitées avec le même soin, une évaluation tenue à part, du code et
des données publiés.

Ce chapitre fixe tout ce qui précède un résultat : les outils et leurs versions, les programmes et leurs
réglages, les tâches, les budgets, les graines aléatoires, le protocole d'évaluation, le format des fichiers,
les statistiques et le plan de chaque expérience. Il décrit ce que le code fait, ligne à ligne quand c'est
utile. Les résultats viennent au
\cref{chap:06-experiences-resultats}. La chaîne logicielle qui relie les essais aux documents est dessinée à
l'annexe~\ref{chap:08-logiciel}.

\section{Un protocole écrit avant le premier entraînement}\label{sec:protocole}

Le protocole des expériences est le fichier \file{sim/SIM\_PROTOCOL.md}. Il fixe les questions, les
algorithmes, les tâches, les budgets, le nombre de graines aléatoires, l'évaluation et les analyses. Les
trois protocoles du projet ont été écrits avant le travail qu'ils encadrent et chaque changement ultérieur y
est consigné et publié dans le dépôt (\file{study/PROTOCOL.md}, \file{sim/SIM\_PROTOCOL.md},
\file{sim/llm/LLM\_PROTOCOL.md}). Dans tout le rapport, un
\emph{essai} désigne un entraînement complet d'un programme sur une tâche avec une graine aléatoire donnée. Le
protocole pose trois règles de conduite :
\begin{itemize}
  \item les graines aléatoires sont fixées avant de lancer, de 0 à 29 ;
  \item aucun essai n'est retiré ni relancé à cause de son score ;
  \item un essai qui plante est relancé une fois avec la même graine aléatoire, puis l'incident est noté.
\end{itemize}

Aucune analyse n'a retiré un essai, changé un réglage ou changé la règle du vainqueur établi
(\cref{sec:regle}).

Les résultats \emph{principaux} sont ceux des expériences E1 à E7 du protocole, à dix graines aléatoires par
couple de programme et de tâche. Les autres analyses les complètent sans les remplacer : B1 et B3, la
consolidation (vingt graines aléatoires, budget égal, réplication) et les contrôles du
\cref{tab:experiences}, qui donne le statut de chacune.

\section{Les outils et leurs versions}\label{sec:outils}

\paragraph{Quatre environnements Python figés}
Quatre environnements Python~3.11.18 ont servi, tous sur un Mac à puce Apple M2 Max (macOS~27.2, arm64).
Chacun est figé par un \texttt{pip freeze}, dans \file{envs/} (\cref{tab:env},
\cref{chap:D-environnements}) :
\begin{itemize}
  \item \file{sim/env} entraîne et réévalue (\file{sim/run.py}, \file{sim/dose.py}) : Stable-Baselines3~2.9.2
    \parencite{raffin2021sb3}, sb3-contrib~2.9.2, SBX~0.28.0 \parencite{raffin2023sbx}, JAX~0.10.2
    \parencite{bradbury2018jax}, Flax~0.12.8, Optax~0.2.8, PyTorch~2.15.0 \parencite{paszke2019pytorch},
    Gymnasium~1.1.1 \parencite{towers2024gymnasium}, MuJoCo~4.0.0 \parencite{todorov2012mujoco} ;
  \item \file{sim/env-an} calcule toutes les analyses : NumPy~2.4.6, SciPy~1.17.1, statsmodels~0.15.0,
    rliable~1.2.0 \parencite{agarwal2021precipice} ;
  \item \file{sim/env-zoo} fait tourner le script officiel de RL Baselines3 Zoo~2.9.2 \parencite{raffin2020zoo}
    pour l'expérience E6, avec les mêmes bibliothèques que \file{sim/env} ;
  \item \file{sim/env-llm} sert au volet des modèles de langue (\cref{chap:07-modeles-de-langue}).
\end{itemize}

\paragraph{Ce que chaque outil apporte}
Gymnasium, successeur d'OpenAI Gym maintenu par la fondation Farama, définit l'interface des tâches :
\texttt{reset} rend une première observation, \texttt{step} applique une action et rend l'observation
suivante, la récompense et deux drapeaux de fin d'épisode. MuJoCo calcule la physique. Il est gratuit et
ouvert depuis 2022, alors qu'il demandait une licence payante en 2018. Stable-Baselines3 fournit PPO écrit en
PyTorch. SBX (\enquote{Stable-Baselines Jax}) réécrit PPO, SAC et CrossQ en JAX avec la même interface.
sb3-contrib fournit TRPO. rliable calcule les moyennes interquartiles, leurs intervalles bootstrap et les
probabilités d'amélioration.

\paragraph{Les versions qui ont tourné}
Tous les entraînements ont tourné avec Stable-Baselines3~2.9.2. Les versions de \file{envs/train.txt} sont
celles qui ont réellement tourné.

\paragraph{Reproduire sur une autre machine}
Chaque entraînement tourne sur un seul fil de calcul (\texttt{OMP\_NUM\_THREADS=1},
\texttt{torch.set\_num\_threads(1)}, options XLA à un fil pour JAX), pour qu'un essai soit déterminé par sa
graine aléatoire sur une machine donnée. Les relances de PPO PyTorch et de PPO JAX l'ont vérifié
(\cref{sec:calcul}). Pour SAC et CrossQ, trois entraînements refaits à l'identique, SAC sur \task{Hopper}
aux graines aléatoires 0 et 1 et CrossQ sur \task{Walker2d} à la graine 0, redonnent les mêmes points
d'évaluation à l'octet près (\file{sim/analysis/c6\_determinism.json}). Tous nos essais ont tourné sur cette seule machine. Sur un
autre processeur, MuJoCo, JAX et PyTorch peuvent enchaîner les calculs en virgule flottante dans un autre ordre et
arrondir autrement : la même graine aléatoire y donnerait alors une autre courbe. Nous ne l'avons pas mesuré. Une
relance ailleurs met à l'épreuve les statistiques du rapport, sans espoir d'en retrouver les courbes au chiffre près.

\section{Les tâches}\label{sec:taches}

Les quatre tâches sont celles de \textcite{henderson2018matters}, dans leur version courante de Gymnasium :
\task{HalfCheetah-v5}, \task{Hopper-v5}, \task{Walker2d-v5} et \task{Swimmer-v5}. Dans chacune, un corps
articulé est simulé par MuJoCo. L'agent observe à chaque pas un vecteur de positions et de vitesses, puis
choisit une action continue : le couple appliqué à chaque articulation. La récompense paie surtout la vitesse
vers l'avant, moins un coût qui croît avec l'effort des moteurs. Pour \task{Hopper} et \task{Walker2d}, elle
ajoute une prime à chaque pas où le corps reste debout. Un épisode s'arrête au plus tard après
\num{1000}~pas.
\begin{itemize}
  \item \task{HalfCheetah} est un guépard plan à six articulations. Il ne tombe jamais : ses épisodes
    durent toujours \num{1000}~pas.
  \item \task{Hopper} est un sauteur à une jambe. L'épisode s'arrête quand il tombe, ce qui rend les scores
    très inégaux d'un essai à l'autre.
  \item \task{Walker2d} est un marcheur plan à deux jambes, arrêté lui aussi à la chute.
  \item \task{Swimmer} est un nageur à trois segments dans un fluide visqueux. Ses scores sont bien plus bas
    que ceux des trois autres tâches.
\end{itemize}

Le \enquote{score} d'un épisode est son retour : la somme des récompenses reçues de son premier à son dernier
pas. Les scores des quatre tâches ne se comparent pas entre eux. Toutes les comparaisons de ce rapport se font
à l'intérieur d'une tâche.

L'expérience E6 emploie aussi \task{Hopper-v4}, la version précédente. Les réglages de PPO que publie RL
Baselines3 Zoo pour Hopper sont inscrits sous ce nom de tâche dans le fichier \file{hyperparams/ppo.yml} du
recueil : ils ont été réglés sur cette version.

\section{Les programmes}\label{sec:programmes}

Nos expériences comparent quatre \emph{programmes}, c'est-à-dire quatre codes prêts à l'emploi, qui réalisent
trois \emph{algorithmes}. Un cinquième programme, TRPO, sert à une seule expérience (B1). Le
\cref{tab:hyperparametres} donne leurs réglages.

\paragraph{PPO, deux fois}
PPO (\emph{Proximal Policy Optimization}, \cite{schulman2017ppo}) est l'algorithme de référence de 2018. Il
sert encore de point de comparaison courant en 2026. C'est aussi l'ancêtre direct de GRPO,
l'algorithme du volet des modèles de langue. PPO est un algorithme \enquote{sur politique} (\emph{on-policy}).
Il collecte un paquet de transitions avec la politique courante, s'en sert pour quelques passes de gradient,
puis le jette. Soit $r_t(\theta) = \pi_\theta(a_t\mid s_t)/\pi_{\theta_{\text{old}}}(a_t\mid s_t)$ le rapport
des probabilités de l'action sous la nouvelle et l'ancienne politique. Soit $\hat A_t$ l'avantage estimé par GAE
(\emph{generalized advantage estimation}) : combien l'action prise a rapporté de plus que ce que la fonction de
valeur prévoyait pour cet état. PPO maximise l'objectif tronqué
\begin{equation}\label{eq:ppo}
  L^{\text{CLIP}}(\theta) = \hat{\mathbb{E}}_t\Bigl[\min\bigl(r_t(\theta)\,\hat A_t,\;
  \operatorname{clip}(r_t(\theta), 1-\epsilon, 1+\epsilon)\,\hat A_t\bigr)\Bigr],
\end{equation}
avec $\epsilon = 0{,}2$, plus un terme d'erreur de la fonction de valeur. Deux programmes réalisent PPO :
\begin{itemize}
  \item \algo{PPO PyTorch}, la classe \texttt{PPO} de Stable-Baselines3 (\texttt{ppo\_sb3} dans le code) ;
  \item \algo{PPO JAX}, la classe \texttt{PPO} de SBX (\texttt{ppo\_sbx} dans le code).
\end{itemize}
Les deux programmes ont les mêmes réglages d'algorithme : chaque ligne du \cref{tab:hyperparametres} a la même
valeur par défaut dans les deux bibliothèques, sauf l'initialisation orthogonale. Les deux bibliothèques
initialisent les poids autrement :
\begin{itemize}
  \item Stable-Baselines3 active par défaut l'option \texttt{ortho\_init}. Chaque matrice de poids est alors
    tirée orthogonale puis multipliée par un gain : $\sqrt2$ pour les couches cachées, 0,01 pour la sortie de
    l'action, 1 pour la sortie de la valeur. Les biais partent de zéro
    (\file{stable\_baselines3/common/policies.py}, lignes 612 à 625).
  \item SBX laisse l'option désactivée et garde l'initialisation de Flax pour toutes les couches : poids tirés
    dans une loi normale tronquée de variance $1/n_{\text{e}}$, où $n_{\text{e}}$ est le nombre d'entrées de
    la couche (loi de LeCun), biais nuls.
\end{itemize}
L'option ne fait pas non plus la même chose dans les deux bibliothèques. Activée dans SBX, elle n'initialise
orthogonalement que la sortie de l'action, au gain 0,01 (\file{sbx/ppo/policies.py}, lignes 81 à 84). Les deux
programmes diffèrent enfin par la bibliothèque de calcul, ses générateurs aléatoires et ses arrondis.
L'expérience E3 mesure ce que ces différences coûtent ensemble, sans les séparer. Un complément, fixé
d'avance, porte les deux programmes à quarante graines aléatoires sur \task{Walker2d} et \task{Swimmer}, puis
échange l'option entre eux pour isoler la part de l'initialisation (\cref{sec:e3-init}).

\paragraph{SAC}
SAC (\emph{Soft Actor-Critic}, \cite{haarnoja2018sac}) sert de référence \enquote{hors politique}
(\emph{off-policy}) en contrôle continu depuis 2019. Il garde toutes les transitions vues dans une mémoire de
rejeu et apprend à chaque pas sur un lot tiré de cette mémoire. Il maximise le retour augmenté de l'entropie de
la politique, $\sum_t \mathbb{E}\bigl[r_t + \alpha\,\mathcal{H}(\pi(\cdot\mid s_t))\bigr]$, avec un
coefficient $\alpha$ réglé automatiquement vers une entropie cible. Deux réseaux critiques estiment la valeur
des actions. Des copies lentes de ces réseaux, dites réseaux cibles, stabilisent l'apprentissage. Le programme
est la classe \texttt{SAC} de SBX (\texttt{sac}).

\paragraph{CrossQ}
CrossQ \parencite{bhatt2024crossq}, paru à ICLR~2024, vise l'efficacité en données, le score atteint par pas
d'interaction avec l'environnement. C'est un descendant de SAC. Il supprime les réseaux cibles, élargit le
critique et insère dans le critique comme dans l'acteur des couches de renormalisation par lots (\emph{batch
renormalization}, \cite{ioffe2017batchrenorm}). Le programme est la classe \texttt{CrossQ} de SBX (\texttt{crossq}). Son critique a deux
couches de \num{1024}~neurones, au lieu des \num{2048} de l'article : SBX fait ce choix pour calculer plus vite.
Nous l'avons gardé. CrossQ emploie aussi Adam avec $\beta_1 = 0{,}5$ au lieu de 0,9, valeur de l'implémentation
d'origine que SBX garde par défaut (\file{sbx/crossq/policies.py}).

\paragraph{La renormalisation par lots de CrossQ, telle que SBX la code}
Une couche de renormalisation centre et réduit chaque activation. La classe \texttt{BatchRenorm} de
\file{sbx/common/jax\_layers.py} le fait en deux phases. Pendant ses \num{100000} premiers passages
d'entraînement, la couche emploie la moyenne et la variance du lot courant, comme une normalisation par lots
ordinaire. Ensuite, elle corrige ces statistiques par des moyennes glissantes d'élan 0,99, que chaque passage met à
jour dès le premier. Le compteur de chaque couche n'avance que lorsque son réseau est entraîné. Le critique est mis à jour à
chaque pas après les 100 premiers, soit près de \num{200000} fois : il passe la moitié de l'entraînement dans la
première phase. L'acteur n'est mis à jour qu'une fois sur trois (\texttt{policy\_delay=3}), soit environ
\num{67000} fois en \num{200000}~pas : ses couches ne quittent jamais la première phase
(\file{sbx/crossq/crossq.py}, où l'acteur et le critique s'appellent l'un l'autre en mode d'évaluation, sans
faire avancer le compteur de l'autre). Nous l'avons mesuré sur un modèle sauvegardé, celui de la graine~0 sur
Hopper (\file{sim/models/crossq/Hopper-v5/seed0.zip}). Le compteur \texttt{steps} des trois couches
\texttt{BatchRenorm} de l'acteur vaut \num{66633}, celui des trois couches du critique \num{199900}, pour un
échauffement de \num{100000} passages. Les autres graines et les autres tâches suivent le même calendrier de
mises à jour, mais nous n'avons lu que ce modèle.

\paragraph{TRPO, pour B1 seulement}
TRPO (\emph{Trust Region Policy Optimization}, \cite{schulman2015trpo}) est l'algorithme de la figure~5 de
\textcite{henderson2018matters}, celle dont l'expérience B1 cherche à expliquer le $t$ de 9
(\cref{sec:plan-b1}). Il maximise l'avantage pondéré par le rapport des
probabilités sous une contrainte de divergence de Kullback-Leibler moyenne entre l'ancienne et la nouvelle
politique, $\hat{\mathbb{E}}_t\bigl[\mathrm{KL}(\pi_{\theta_{\text{old}}}\Vert\pi_\theta)\bigr] \le \delta$,
résolue par gradient conjugué puis recherche linéaire. Le programme est la classe \texttt{TRPO} de sb3-contrib
(\texttt{trpo}), avec toutes ses valeurs par défaut.

\paragraph{Pourquoi les valeurs par défaut}
Les quatre programmes de la grille principale gardent les valeurs par défaut de leur bibliothèque. Aucun n'a
été réglé. Ce choix applique à la lettre la demande de Pineau sur les références (faille F7, des références moins bien réglées que la méthode) : chaque
programme
reçoit le même soin de réglage, en l'occurrence aucun. Cette égalité a un coût : un programme bien servi par ses
valeurs par défaut est avantagé. Le classement de E2 vaut pour ces réglages. L'expérience E6 mesure, sur un cas, ce que change un
réglage publié. La seule commande qui construit un programme de la grille, dans \file{sim/run.py}, est
\texttt{cls("MlpPolicy", TASK, seed=SEED, verbose=0)}, sans autre argument. Le \cref{tab:hyperparametres}
relève les réglages tels qu'ils ont servi : \file{sim/analysis/hyperparams.py} reconstruit chaque programme comme
\file{sim/run.py}, lit les attributs de l'algorithme, de la politique et de l'optimiseur, puis les compare à ceux
d'un modèle enregistré.

\gentable{hyperparametres}

\paragraph{Les réglages publiés du Zoo, pour E6}
L'expérience E6 oppose PPO PyTorch avec ses valeurs par défaut à PPO PyTorch avec les réglages que RL
Baselines3 Zoo publie pour \task{Hopper-v4}. \file{sim/run.py} les recopie dans la constante
\texttt{ZOO\_PPO\_HOPPER}, identiques à ceux de \file{rl\_zoo3/hyperparams/ppo.yml} dans la version~2.9.2 :
\begin{itemize}
  \item paquets de 512~pas, lots de 32, 5~passes, $\gamma = 0{,}999$, $\lambda = 0{,}99$ pour GAE ;
  \item taux d'apprentissage \num{9.80828e-5}, coefficient d'entropie \num{0.00229519}, coefficient de
    valeur \num{0.835671}, norme de gradient bornée à 0,7, troncature 0,2 ;
  \item réseaux de deux couches de 256~neurones pour la politique et la valeur, activation ReLU, sans
    initialisation orthogonale, logarithme de l'écart type initial des actions à $-2$ ;
  \item normalisation des observations et des récompenses par une moyenne et un écart type glissants
    (\texttt{VecNormalize}).
\end{itemize}
La normalisation des récompenses divise chaque récompense par l'écart type d'un retour actualisé glissant.
Dans la recette du Zoo, son facteur d'actualisation est celui de l'algorithme, 0,999. Notre premier passage
(\texttt{ppo\_sb3\_zoo}, option \texttt{-{}-zoo}) a laissé par erreur la valeur par défaut de
\texttt{VecNormalize}, 0,99. Nous avons corrigé cette erreur par un second passage (\texttt{ppo\_sb3\_zoofix},
option \texttt{-{}-zoofix}), qui passe 0,999. Les deux passages sont conservés. Le diagnostic de E6 les
compare : la correction ne change pas le résultat de façon détectable (\cref{sec:e6}).

\section{Les budgets}\label{sec:budgets}

\paragraph{Ce que fixe le protocole}
PPO PyTorch, PPO JAX et TRPO s'entraînent un million de pas. SAC et CrossQ s'entraînent \num{200000}~pas. Un
pas est une transition de l'environnement : une observation, une action, une récompense.

\paragraph{Pourquoi deux budgets}
Les deux familles d'algorithmes ne se comparent pas au même régime. Un algorithme sur politique jette ses
données après chaque mise à jour : il lui faut beaucoup de pas. Un million est l'ordre de grandeur de 2018
pour ces tâches. Un algorithme hors politique réutilise ses données à chaque pas : les articles de 2024 à 2026 comparent SAC et
CrossQ dans un régime \enquote{peu de données}. Le budget de \num{200000}~pas tient le calcul sur une
machine partagée : ces deux algorithmes y prennent 19~ms par pas.

\paragraph{Le nombre exact de pas}
\file{sim/run.py} entraîne par blocs de \num{50000}~pas et évalue après chaque bloc. L'appel est
\texttt{model.learn(50\_000, reset\_num\_timesteps=False)}, répété jusqu'au budget. PPO et TRPO collectent leurs transitions par
paquets de \num{2048}~pas et ne s'arrêtent qu'à la fin d'un paquet. Chaque bloc fait donc \num{51200}~pas.
Les modèles enregistrés de PPO PyTorch, PPO JAX et TRPO comptent \num{1024000}~pas, alors que la colonne
\texttt{step} des résultats écrit \num{1000000}. Ceux de E6 réglés comme le Zoo, à paquets de 512~pas, en
comptent \num{1003520}. SAC et CrossQ font exactement \num{200000}~pas.

\paragraph{Le budget égal}
La comparaison de E2 se fait à budgets inégaux. Le rapport le dit à chaque fois. Une lecture à budget égal la
complète : PPO est lu à \num{200000}~pas sur sa courbe d'évaluation, au même point que SAC et CrossQ.

\section{Les graines aléatoires}\label{sec:graines}

\paragraph{Ce que fixe une graine aléatoire}
Une graine aléatoire est le nombre de départ dont l'ordinateur tire tout le hasard d'un entraînement :
l'initialisation des réseaux, les positions initiales du corps simulé, le bruit des actions explorées, les lots
tirés dans la mémoire de rejeu. \file{sim/run.py} la passe au constructeur du programme
(\texttt{seed=SEED}), qui en dérive les générateurs de NumPy, de PyTorch ou de JAX et celui de
l'environnement d'entraînement. Pour PPO PyTorch et PPO JAX, une relance avec la même graine aléatoire sur la
même machine a redonné les mêmes points d'évaluation (\cref{sec:calcul}). Pour SAC et CrossQ, trois
entraînements refaits à l'identique les redonnent à l'octet près (\cref{sec:outils}).

\paragraph{Graines aléatoires d'entraînement}
Les graines aléatoires d'entraînement sont les entiers 0 à 29, fixés avant le premier entraînement :
\begin{itemize}
  \item 0 à 9 pour chaque couple de programme et de tâche, les résultats principaux ;
  \item 10 à 19 pour la consolidation, les vingt graines aléatoires ;
  \item 0 à 29 pour PPO PyTorch sur \task{HalfCheetah} (E1) et pour TRPO sur \task{HalfCheetah} (B1).
\end{itemize}
Les graines aléatoires 10 à 19 de PPO PyTorch sur \task{HalfCheetah}, entraînées pour E1 avec le même code, servent telles
quelles à la consolidation.

\paragraph{Graines aléatoires d'évaluation}
L'évaluation emploie les graines aléatoires \num{10000} à \num{10009}, une par épisode d'évaluation. Elles ne
servent jamais à l'entraînement, qui emploie les graines 0 à 29. Les mêmes dix graines aléatoires servent à chaque point
d'évaluation, à chaque essai et à chaque physique : deux politiques sont toujours évaluées sur les mêmes dix
positions de départ.

\paragraph{Graines aléatoires des analyses}
Les tirages des analyses ont leurs propres graines aléatoires, fixées dans chaque script : 2018 pour les tirages de E1 et
B1 (\texttt{np.random.default\_rng(2018)}), 2026 pour E5 et pour tous les bootstraps de rliable
(\cref{sec:reproductibilite-bootstrap}). Deux exécutions d'une analyse rendent un fichier identique à l'octet
près.

\paragraph{Chaque script nomme ses graines aléatoires}
Six de nos scripts d'analyse lisaient tous les fichiers présents. Avec les graines aléatoires 10 à 19 de la
consolidation, E6 serait passé de dix à vingt graines aléatoires sans le dire. Nous avons corrigé ce défaut :
chaque script nomme ses graines aléatoires et s'arrête s'il en manque une. Les résultats publiés n'ont pas
bougé.

\section{L'évaluation, tenue à part}\label{sec:evaluation}

\paragraph{Le protocole d'évaluation}
Après chaque bloc d'entraînement, la fonction \texttt{evaluate} de \file{sim/run.py} joue dix épisodes. Pour
l'épisode $i$, de 0 à 9 :
\begin{enumerate}
  \item un environnement neuf est créé (\texttt{gym.make(TASK)}), distinct de celui de l'entraînement ;
  \item il est remis à zéro avec la graine aléatoire $\num{10000}+i$ ;
  \item la politique choisit chaque action sans hasard (\texttt{deterministic=True}) : l'action moyenne pour
    PPO, SAC et CrossQ, au lieu d'un tirage dans la loi de la politique ;
  \item l'épisode court jusqu'à la chute ou jusqu'à \num{1000}~pas, puis son retour est noté.
\end{enumerate}
La fonction rend la moyenne et l'écart type des dix retours (\texttt{np.std}, sans correction de Bessel). Pour
les programmes réglés comme le Zoo, chaque observation passe par la normalisation de l'entraînement, dont les
statistiques ne bougent pas pendant l'évaluation. Une politique déterministe évaluée sur dix graines aléatoires fixes
donne toujours le même score. La variation entre les dix épisodes ne vient que des positions de départ.

\paragraph{Le score d'un essai}
Le score d'un essai est la moyenne des dix retours à la dernière évaluation, au budget de son programme : la
colonne \texttt{eval\_mean} de la dernière ligne de son fichier. C'est la mesure principale fixée par le
protocole. L'expérience B3 en essaie deux autres.

\paragraph{Une physique modifiée, à la fin}
À la fin de l'entraînement, la politique finale joue dix épisodes de plus, sur les mêmes graines aléatoires
d'évaluation, dans une physique modifiée. Le code modifie le modèle MuJoCo de l'environnement avant la remise
à zéro :
\begin{itemize}
  \item la masse de chaque corps est multipliée par 1,2 (\texttt{body\_mass[:] *= 1.2}) ;
  \item le coefficient de frottement de glissement de chaque géométrie, sol compris, est multiplié par 0,8
    (\texttt{geom\_friction[:, 0] *= 0.8}).
\end{itemize}
La moyenne des dix retours s'écrit dans la colonne \texttt{perturbed\_mean} de la dernière ligne. Elle sert à
l'expérience E7.

\paragraph{La physique modifiée par degrés}
La consolidation réévalue chaque politique finale dans d'autres physiques, sans aucun entraînement
(\file{sim/dose.py}). Une seule grandeur change à la fois :
\begin{itemize}
  \item les masses, multipliées par 0,7, 0,8, 0,9, 1,0, 1,1, 1,2 et 1,3, soit de 30~\% plus légères à
    30~\% plus lourdes ;
  \item les frottements de glissement, multipliés par 0,6, 0,7, 0,8, 0,9, 1,1 et 1,2, soit de 40~\% plus
    faibles à 20~\% plus forts.
\end{itemize}
Chaque physique est jouée sur les dix mêmes graines aléatoires d'évaluation, avec le même code d'évaluation. Le facteur
1,0 sert de contrôle : une réévaluation correcte y redonne le score final écrit par \file{sim/run.py}. Le résultat de
chaque politique est un fichier \file{results/dose/}\allowbreak\file{<programme>/}\allowbreak\file{<tâche>/}\allowbreak\file{seed<k>.json}. Les politiques de PPO
de la grille principale n'ont pas de normalisation des observations : \file{sim/dose.py} les évalue telles
quelles.

\section{Le format des essais}\label{sec:format}

\paragraph{Un fichier CSV par essai}
Chaque essai écrit un fichier
\file{results/}\allowbreak\file{<programme>/}\allowbreak\file{<tâche>/}\allowbreak\file{seed<k>.csv}, réécrit après chaque évaluation.
Ses colonnes sont
\texttt{step,}\allowbreak\texttt{eval\_mean,}\allowbreak\texttt{eval\_std,}\allowbreak\texttt{perturbed\_mean,}\allowbreak\texttt{wall} :
\begin{itemize}
  \item \texttt{step}, le nombre de pas annoncé à cette évaluation, de \num{50000} en \num{50000} ;
  \item \texttt{eval\_mean} et \texttt{eval\_std}, la moyenne et l'écart type des retours des dix épisodes
    d'évaluation ;
  \item \texttt{perturbed\_mean}, la moyenne des retours dans la physique modifiée, vide sauf à la dernière
    ligne ;
  \item \texttt{wall}, les secondes écoulées depuis le début de l'essai.
\end{itemize}
Une courbe de PPO ou de TRPO a vingt lignes, une courbe de SAC ou de CrossQ en a quatre. Le module
\file{src/pineau\_rl/runs.py} lit ce format. Sa fonction \texttt{read\_run} vérifie les colonnes requises et
la croissance stricte des pas. Elle tient pour inachevé un essai dont la dernière ligne n'est pas au budget du
programme. Elle coupe aussi, sur demande, une courbe à un pas donné (lecture à budget égal). La fonction
\texttt{write\_run} écrit ce format depuis n'importe quelle bibliothèque.

\paragraph{Les politiques finales}
La politique finale de chaque essai est enregistrée dans
\file{models/}\allowbreak\file{<programme>/}\allowbreak\file{<tâche>/}\allowbreak\file{seed<k>.zip}, avec les
statistiques de normalisation à côté pour les programmes réglés comme le Zoo (\file{.vecnorm.pkl}). Le fichier
\file{sim/models/MANIFEST.sha256} donne l'empreinte SHA-256 de 250~politiques, prises sur
la machine qui les a entraînées. Les politiques de PPO, de TRPO et de E6 sont dans le dépôt. Celles de SAC et
de CrossQ, trop lourdes pour git, sont attachées à la publication \texttt{policies-v1} du dépôt sur GitHub.
Le script \file{sim/fetch\_models.py} les télécharge une à une, espacées d'une pause aléatoire de 1,5 à
4~secondes, vérifie chaque archive puis chaque politique contre le manifeste. L'option
\texttt{-{}-check} vérifie seulement les politiques présentes.

\paragraph{Les journaux}
Un journal accompagne chaque lancement. Pour la consolidation, chaque tentative garde le sien
(\file{logs/<tâche>.attempt<n>.log}, jamais écrasé) et ajoute une ligne au registre
\file{logs/attempts.jsonl} : tâche, numéro de tentative, début, fin, code de sortie, empreinte SHA-256 du
résultat. Pour les entraînements antérieurs, le registre \file{logs/attempts\_reconstructed.csv} reconstruit une ligne
par tentative à partir des dates des journaux, des sorties des lanceurs et de l'historique git
(\cref{sec:calcul}).

\section{Les statistiques}\label{sec:statistiques}

Cette section définit chaque statistique employée, dans l'ordre où le \cref{chap:06-experiences-resultats} en
a besoin. Les fonctions sont dans
\file{src/pineau\_rl/stats.py} et \file{src/pineau\_rl/}\allowbreak\file{bootstrap.py}. Les statistiques des parts
du sondage, l'intervalle de Wilson d'une part, celui de Newcombe d'une différence de deux parts et le test exact
de Fisher, sont définies avec leurs formules à la \cref{sec:sm-stats}. Les expériences ne les emploient pas.
Leurs parts à elles, comme un taux de faux positifs, sont des fréquences sur $N$ tirages simulés : leur erreur
de Monte-Carlo vaut $\sqrt{\hat p(1-\hat p)/N}$ quand les tirages sont indépendants.

\subsection{La moyenne interquartile}\label{sec:iqm}
La moyenne interquartile (IQM, \emph{interquartile mean}) écarte le quart le plus bas et le quart le plus haut
des scores, puis fait la moyenne du reste \parencite{agarwal2021precipice}. Elle résiste mieux qu'une moyenne à
un essai hors norme et reste plus précise qu'une médiane. rliable la calcule par la moyenne tronquée de SciPy à
25~\% de chaque côté (\texttt{trim\_mean(x, 0.25)}), qui retire $\lfloor n/4\rfloor$ scores à chaque extrémité.
Avec les scores triés $x_{(1)}\le\dots\le x_{(n)}$ et $g=\lfloor n/4\rfloor$ :
\begin{equation}\label{eq:iqm}
  \mathrm{IQM}(x) = \frac{1}{n-2g}\sum_{i=g+1}^{n-g} x_{(i)} .
\end{equation}
Avec dix graines aléatoires, $g=2$ : l'IQM est la moyenne des six scores du milieu. Avec vingt, $g=5$ et dix
scores restent. Avec trente, $g=7$ et seize restent.

\subsection{L'intervalle bootstrap stratifié}\label{sec:bootstrap}
L'incertitude d'une IQM se mesure par bootstrap \parencite{efron1979bootstrap}. On tire avec remise $n$
scores parmi les $n$ observés, on recalcule l'IQM, puis on répète $B$ fois. L'intervalle à 95~\% est formé des
centiles 2,5 et 97,5 des $B$ valeurs obtenues (méthode des centiles, la méthode par défaut de rliable) :
\begin{equation}\label{eq:bootstrap}
  \bigl[\,q_{0{,}025}\bigl(\mathrm{IQM}(x^{*1}),\dots,\mathrm{IQM}(x^{*B})\bigr),\;
  q_{0{,}975}\bigl(\mathrm{IQM}(x^{*1}),\dots,\mathrm{IQM}(x^{*B})\bigr)\,\bigr],
  \qquad B = \num{2000}.
\end{equation}
rliable appelle ce bootstrap \enquote{stratifié} : avec plusieurs tâches, il rééchantillonne les essais à
l'intérieur de chaque tâche, sans mêler les tâches. Chacun de nos appels traite une seule tâche.
La stratification n'y change rien : le bootstrap revient à rééchantillonner les essais d'un programme.

\subsection{La probabilité d'amélioration}\label{sec:poi}
La probabilité d'amélioration d'un programme $X$ sur un programme $Y$ est la probabilité qu'un essai de $X$,
tiré au hasard, batte un essai de $Y$, tiré au hasard, les égalités comptant pour moitié
\parencite{agarwal2021precipice}. Avec $n_x$ scores $x_i$ et $n_y$ scores $y_j$ :
\begin{equation}\label{eq:poi}
  P(X>Y) = \frac{1}{n_x n_y}\sum_{i=1}^{n_x}\sum_{j=1}^{n_y}
  \Bigl(\mathbf{1}[x_i>y_j] + \tfrac12\,\mathbf{1}[x_i=y_j]\Bigr) = \frac{U}{n_x n_y},
\end{equation}
où $U$ est la statistique de Mann-Whitney \parencite{mann1947test}. rliable la calcule ainsi
(\texttt{scipy.stats.mannwhitneyu}). Elle vaut 0,5 quand les deux programmes se valent. Son intervalle à
95~\% vient du même bootstrap que l'IQM : les essais de $X$ et ceux de $Y$ sont rééchantillonnés chacun de leur
côté, $B=\num{2000}$ fois (classe \texttt{StratifiedIndependentBootstrap} de rliable), puis l'on prend les
centiles 2,5 et 97,5. À un niveau $1-\alpha$, on prend les centiles $100\,\alpha/2$ et $100\,(1-\alpha/2)$.

\subsection{La règle du vainqueur}\label{sec:regle}
Le protocole fixe une règle unique pour déclarer qu'un programme l'emporte sur un autre. Nous l'appelons
\emph{la règle du vainqueur établi}. Une différence est établie quand l'une de deux conditions tient :
\begin{itemize}
  \item les intervalles à 95~\% des deux IQM ne se chevauchent pas ;
  \item l'intervalle à 95~\% de la probabilité d'amélioration exclut 0,5.
\end{itemize}
Sinon, aucun vainqueur n'est déclaré, quelle que soit la différence entre les deux IQM. E3 applique les deux conditions
(\file{e3.py}). E2, E2 avec Bonferroni et B3 emploient la seconde : une paire y est dite \enquote{tranchée}
quand l'intervalle de sa probabilité d'amélioration exclut 0,5. E7 et la physique par degrés en emploient
l'analogue pour un rapport : une perte est établie quand l'intervalle de l'IQM du rapport reste sous 1.

\paragraph{Le niveau réel de la règle}
Deux conditions réunies par un \enquote{ou} n'ont pas de niveau garanti. Un intervalle bootstrap percentile
calculé sur dix essais n'a qu'un niveau approché. Nous mesurons le niveau de la règle là où aucune différence n'existe : deux groupes disjoints tirés parmi les trente essais d'un même
programme (\file{sim/analysis/rule\_calibration.py}, \num{2000} tirages, \num{2000} rééchantillons par
intervalle). Une règle exacte au niveau de 95~\% y déclarerait un vainqueur dans 5~\% des tirages. Le
\cref{tab:regle-tranche} donne le taux réel : 5,4~\% pour PPO PyTorch et 6,2~\% pour TRPO à dix essais contre
dix, 9,9~\% et 9,1~\% à cinq contre cinq. Sur \num{2000}~tirages indépendants, l'erreur de Monte-Carlo d'une part
voisine de 5~\% vaudrait environ un demi-point. Nos tirages, pris dans les trente mêmes essais, ne sont pas
indépendants : ce demi-point ne donne qu'un ordre de grandeur. À cette réserve près, 5,4~\% ne se distingue pas
de 5~\%, alors que les taux à cinq contre cinq s'en éloignent sans ambiguïté. La condition sur la
probabilité d'amélioration fait presque tout le travail. Les intervalles des IQM ne sont disjoints que dans
0,6~\% et 0,9~\% des tirages à dix contre dix. La règle est donc un peu libérale à dix graines aléatoires et
nettement libérale à cinq. Un verdict dont l'intervalle frôle 0,5 se lit avec cette réserve, comme celui de
PPO JAX contre CrossQ sur \task{Hopper} à dix graines aléatoires, $[0{,}52\,;\,0{,}94]$.

\gentable{regle-tranche}

\subsection{Les tests t de Student et de Welch}\label{sec:welch}
Le test t compare deux moyennes. Avec deux groupes de tailles $n_a$ et $n_b$, de moyennes $\bar a$ et $\bar b$
et de variances empiriques $s_a^2$ et $s_b^2$ (corrigées de Bessel), le test de Student
\parencite{student1908probable} suppose les deux variances égales et calcule
\begin{equation}\label{eq:student}
  t = \frac{\bar a-\bar b}{s_p\sqrt{1/n_a+1/n_b}},\qquad
  s_p^2=\frac{(n_a-1)s_a^2+(n_b-1)s_b^2}{n_a+n_b-2},
\end{equation}
avec $n_a+n_b-2$ degrés de liberté. Le test de Welch \parencite{welch1947generalization} ne suppose pas les
variances égales :
\begin{equation}\label{eq:welch}
  t = \frac{\bar a-\bar b}{\sqrt{s_a^2/n_a+s_b^2/n_b}},\qquad
  \nu = \frac{\bigl(s_a^2/n_a+s_b^2/n_b\bigr)^2}{\dfrac{(s_a^2/n_a)^2}{n_a-1}+\dfrac{(s_b^2/n_b)^2}{n_b-1}} .
\end{equation}
La valeur $p$ bilatérale est $2\,\bigl(1-F_\nu(|t|)\bigr)$, où $F_\nu$ est la fonction de répartition de la
loi de Student à $\nu$ degrés de liberté. Deux groupes de cinq donnent $\nu\ge 4$ : $\nu$ vaut au moins le plus
petit des $n-1$. Tous nos tests t sont des tests de Welch (\texttt{scipy.stats.ttest\_ind}, option
\texttt{equal\_var=False}). C'est le test que l'expérience E1 met en cause, comme le prévoyait le protocole.

\subsection{Trois lectures du t moyen de 2018}\label{sec:t-moyen}
La légende de la figure~5 de \textcite{henderson2018matters} parle du \enquote{test t moyen à deux
échantillons sur toute la distribution d'entraînement}. Elle compare deux groupes de courbes d'apprentissage entières.
Soit $C^a_{s,k}$ le score d'évaluation de l'essai $s$ du groupe $a$ au point
d'évaluation $k$, pour $k=1,\dots,K$ ($K=20$ pour PPO et TRPO). Nos analyses calculent trois lectures de cette
légende (\texttt{welch\_curve} et \texttt{welch\_pooled}).
\begin{itemize}
  \item \emph{Le t moyen.} Un test de Welch à chaque point, sur les cinq scores de chaque groupe, donne $t_k$ et
    $\nu_k$. On rapporte $\bar t = \frac1K\sum_k t_k$, converti en $p$ avec la moyenne des $\nu_k$.
  \item \emph{Le p moyen.} Le même calcul, résumé par $\bar p = \frac1K\sum_k p_k$.
  \item \emph{Les points empilés.} Les $5\times K$ points de chaque groupe sont pris comme autant
    d'observations indépendantes, puis comparés par un seul test de Welch, à environ 198~degrés de liberté.
\end{itemize}

\subsection{La statistique du script de Henderson}\label{sec:script-henderson}
Peter Henderson a publié le 17~juin~2017 le script \file{from\_same\_dist.py}
\parencite{henderson2017samedist}, compagnon de \textcite{islam2017reproducibility}. Le script fait un autre
calcul que les trois lectures précédentes. La fonction \texttt{henderson\_script} le reproduit sur nos
fichiers. Soit $M_{s,k}$ la colonne \texttt{eval\_mean} et $S_{s,k}$ la colonne \texttt{eval\_std} de
l'essai $s$ au point $k$. Soit $A$ et $B$ les deux groupes de cinq essais.
\begin{enumerate}
  \item Le script fait la moyenne des courbes de chaque groupe, colonne par colonne :
    $\bar M^A_k=\frac15\sum_{s\in A}M_{s,k}$ et $\bar S^A_k=\frac15\sum_{s\in A}S_{s,k}$.
  \item À chaque point $k$, il applique un test de Welch construit à partir de ces seules statistiques
    (\texttt{scipy.stats.ttest\_ind\_from\_stats}), avec $n$ épisodes par groupe :
    \begin{equation}\label{eq:script}
      t_k=\frac{\bar M^A_k-\bar M^B_k}{\sqrt{\bigl((\bar S^A_k)^2+(\bar S^B_k)^2\bigr)/n}} .
    \end{equation}
  \item Il rapporte la moyenne des $t_k$ et la moyenne des $p_k$.
\end{enumerate}
Le dénominateur de l'\cref{eq:script} contient l'écart type \emph{entre les épisodes} d'un même essai. La
variabilité \emph{entre les graines aléatoires}, celle que la figure de 2018 devait mesurer, n'y entre jamais.
Les cinq graines aléatoires d'un groupe se fondent en une seule courbe moyenne. Pour nos essais, $n=10$, le nombre
d'épisodes d'évaluation. Une analyse de sensibilité prend $n=5$, le nombre de trajectoires par itération des essais de 2018
(\num{5000}~pas par itération, épisodes de \num{1000}~pas). Notre $S_{s,k}$ est un écart type sans correction
de Bessel, passé tel quel au test, comme le script passe ses propres colonnes.

\paragraph{De combien le script gonfle le t}
Au dernier point, le $t$ du script se compare à celui d'un test de Welch sur les cinq scores
\texttt{eval\_mean} de chaque groupe. Les deux ont le même numérateur, la différence des moyennes des deux groupes.
Soit $s_A^2$ et $s_B^2$ les variances des scores entre graines aléatoires dans chaque groupe. Soit $\bar S_A$
et $\bar S_B$ les écarts types entre épisodes moyennés sur chaque groupe. Posons $\sigma^2=(s_A^2+s_B^2)/2$ et
$S^2=(\bar S_A^2+\bar S_B^2)/2$. Le test sur les graines a pour erreur type $\sigma\sqrt{2/5}$, le script
$S\sqrt{2/n}$, d'où l'identité exacte
\begin{equation}\label{eq:facteur-b1}
  \frac{t_{\text{script}}}{t_{\text{graines}}}=\frac{\sigma\sqrt{2/5}}{S\sqrt{2/n}}
  =\frac{\sqrt{n/5}}{\rho},\qquad \rho=\frac{S}{\sigma}.
\end{equation}
Avec nos dix épisodes, le facteur vaut $\sqrt2/\rho$. Il ne vaut $1/\rho$ que pour l'analyse de sensibilité à
$n=5$. Nous vérifions cette identité sur les tirages de \file{b1\_script.py}
(\file{sim/analysis/b1\_factor.py}) et rapporte la médiane du facteur, $\rho$ étant recalculé dans chaque tirage.

\subsection{Corrections pour comparaisons multiples}\label{sec:multiples}
Quand on fait $m$ tests, la probabilité qu'au moins un se trompe dépasse le seuil de chacun. Deux corrections
servent.
\begin{itemize}
  \item \emph{Bonferroni} \parencite{dunn1961multiple} élève chaque intervalle au niveau $1-0{,}05/m$. Les
    $m$ intervalles tiennent alors ensemble au niveau de 95~\% au moins. E2 compte $m=24$ comparaisons
    (6~paires sur 4~tâches) : chaque intervalle passe au niveau $1-0{,}05/24$, soit 99,79~\%.
  \item \emph{Holm} \parencite{holm1979simple} ajuste des valeurs $p$ triées
    $p_{(1)}\le\dots\le p_{(m)}$ en
    $\tilde p_{(i)}=\min\bigl(1,\max_{j\le i}(m-j+1)\,p_{(j)}\bigr)$. Une hypothèse est rejetée quand
    $\tilde p_{(i)}<0{,}05$. La probabilité d'au moins un faux rejet dans la famille reste sous 5~\%, sans
    hypothèse sur la dépendance des tests. Holm sert à E3 sur ses quatre tâches et à E6,
    séparément sur les trois comparaisons de son diagnostic et sur les deux de son pipeline officiel.
\end{itemize}
Une même famille peut recevoir les deux corrections : E3 corrigé donne ses valeurs $p$ de Mann-Whitney ajustées
par Holm et ses intervalles de probabilité d'amélioration au niveau de Bonferroni, $1-0{,}05/4$, soit
98,75~\%.

\subsection{Autres tests}\label{sec:autres-tests}
Le test de Shapiro-Wilk \parencite{shapiro1965analysis} mesure l'ajustement d'un échantillon à une loi normale.
E1 l'applique aux trente scores de PPO PyTorch.

\paragraph{Le test de Mann-Whitney}
Le test de Mann-Whitney \parencite{mann1947test} compare deux groupes par leurs rangs. Sa statistique est le $U$
de l'\cref{eq:poi}. Sous l'hypothèse que les deux groupes viennent de la même loi continue, $U$ a pour moyenne
$n_xn_y/2$ et pour variance $n_xn_y(n_x+n_y+1)/12$, aux égalités près. SciPy calcule la valeur $p$ bilatérale
exacte quand un groupe compte au plus huit scores sans égalité, sinon par l'approximation normale, corrigée des
égalités et de la continuité (\texttt{scipy.stats.mannwhitneyu}, méthode \texttt{auto}). Le diagnostic de E6 le
rapporte à côté du test de Welch. E3 corrigé et le test de E7 l'emploient pour leurs valeurs $p$.

\paragraph{La différence de Hodges-Lehmann}
Une probabilité d'amélioration dit qui gagne, pas de combien. La différence en points de score que nous mettons à
côté est l'estimateur de \textcite{hodges1963estimates}, la médiane des $n_xn_y$ différences deux à deux :
\begin{equation}\label{eq:hl}
  \hat\Delta=\operatorname{med}\{\,x_i-y_j : 1\le i\le n_x,\ 1\le j\le n_y\,\}.
\end{equation}
Son intervalle à 95~\% est un bootstrap percentile, chaque groupe rééchantillonné de son côté, \num{5000} fois,
avec \texttt{default\_rng(2026)} (\file{sim/analysis/effect\_sizes.py}, \file{e3\_holm.py}).

\paragraph{Puissance et plus petit effet détectable}
La puissance d'un test est la probabilité qu'il rejette l'hypothèse nulle quand un effet donné existe.
\file{e3\_holm.py} l'estime par rééchantillonnage : \num{2000} fois, il tire avec remise $n$ essais de chaque
programme parmi les vingt observés, applique le test de Mann-Whitney, puis compte la part des tirages où
$p<0{,}05/4$. Il en déduit aussi la plus petite probabilité d'amélioration qu'un test à $n$ contre $n$ détecte
avec 80~\% de puissance, par l'approximation normale de $U$ :
\begin{equation}\label{eq:mde}
  P_{\min}(n)=\frac12+\bigl(z_{1-\alpha/2}+z_{0{,}8}\bigr)\sqrt{\frac{2n+1}{12\,n^2}},\qquad
  \alpha=0{,}05/4,
\end{equation}
où $z_q$ est le quantile d'ordre $q$ de la loi normale centrée réduite et $(2n+1)/(12n^2)$ la variance de
$U/n^2$ sous l'hypothèse nulle.

\paragraph{Le nombre de points indépendants d'une courbe}
Les vingt points d'une courbe d'apprentissage ne sont pas vingt mesures indépendantes : un essai au-dessus de la
courbe moyenne à un point le reste au point suivant. Nous mesurons cette dépendance par l'autocorrélation
d'ordre 1, $r$, de la différence entre chaque courbe et la courbe moyenne, puis en tire le nombre effectif de points
indépendants d'une courbe de 20 points, $20\,(1-r)/(1+r)$, la formule approchée, pour une longue série, d'un
processus autorégressif d'ordre 1
(\file{sim/analysis/curves.py}). La lecture empilée de B1 compte ces vingt points comme vingt observations.

\subsection{Des intervalles reproductibles}\label{sec:reproductibilite-bootstrap}
rliable~1.2.0 accepte un argument \texttt{random\_state} mais tire ses rééchantillons dans le générateur global
de NumPy (\texttt{np.random.choice}). Six de nos scripts ne redonnaient donc pas les mêmes intervalles d'une
exécution à l'autre. Nous avons corrigé ce défaut : chaque script fixe ce générateur au départ
(\texttt{np.random.seed(2026)}), puis appelle rliable dans un ordre fixe. Comme tout appel consomme le
générateur global, l'intervalle d'une comparaison dépend des appels qui la précèdent dans le même script. La correction
n'a changé aucune estimation ponctuelle ni aucun verdict. Les bornes ont bougé de l'ordre du bruit de
Monte-Carlo, au plus 16~points sur une IQM de \num{2423} pour E6.

\paragraph{Une variante pour le diagnostic de E6}
Les scripts \file{e6\_diag.py} et \file{e6\_official.py} calculent les intervalles des IQM avec
$B=\num{5000}$ rééchantillons. Leur probabilité d'amélioration a un bootstrap à part
(\texttt{probability\_of\_improvement\_bootstrap}) : \num{5000} tirages, les deux groupes rééchantillonnés
l'un après l'autre avec un générateur de graine aléatoire 2018, sans rliable. Le même score peut donc recevoir deux
intervalles un peu différents selon le script qui le calcule : l'IQM de PPO PyTorch par défaut sur
\task{Hopper-v5} a un intervalle dans \file{e6.json} et un autre dans \file{e6\_diag.json}. Le rapport cite
celui du diagnostic, comme \file{docs/conclusions.md}.

\section{Les expériences, une à une}\label{sec:plans}

Le \cref{tab:experiences} résume les plans. Il ne contient que des faits de conception : la question, le plan,
le nombre d'essais et le statut de l'analyse.

{\small
\begin{longtable}{@{}>{\raggedright\hyphenpenalty=10000\exhyphenpenalty=10000\arraybackslash}p{2.1cm}
  >{\raggedright\arraybackslash}p{3.4cm}>{\raggedright\arraybackslash}p{3.85cm}
  >{\raggedright\arraybackslash}p{2.0cm}>{\raggedright\arraybackslash}p{2.5cm}@{}}
\caption{Les expériences : la question posée, le plan, les essais employés et le statut de l'analyse.
\enquote{Principale} désigne une expérience du protocole, \enquote{complémentaire} une analyse qui la
complète sans la remplacer, \enquote{exploratoire} un contrôle de plus, dont le résultat se lit comme une
indication. Aucun chiffre de cette
table n'est un résultat.}\label{tab:experiences}\\
\toprule
Expérience & Question & Plan & Essais & Statut \\
\midrule
\endfirsthead
\toprule
Expérience & Question & Plan & Essais & Statut \\
\midrule
\endhead
\bottomrule
\endfoot
E1 & Deux groupes d'essais du même programme passent-ils pour deux algorithmes ? &
PPO PyTorch sur \task{HalfCheetah}, 30 graines aléatoires. \num{10000} tirages de deux groupes, test de Welch
sur les scores finaux, au hasard puis triés. & 30 essais & principale \\
E1, seconde statistique & Le t moyen sur la courbe de 2018 se retrouve-t-il ? &
Les mêmes 30 courbes, test de Welch à chacun des 20 points, t moyen et p moyen. & les 30 de E1 &
complémentaire \\
B1, TRPO & La dispersion de TRPO explique-t-elle le t de 9 ? &
TRPO sur \task{HalfCheetah}, 30 graines aléatoires, trois lectures du t moyen. & 30 essais &
complémentaire \\
B1, script & Le calcul de Henderson explique-t-il le t de 9 ? &
Le script de 2017 appliqué aux courbes de TRPO et de PPO, mêmes tirages. & les 30 de E1 et les 30 de TRPO &
complémentaire \\
E2 & Le programme en tête change-t-il de tâche en tâche ? &
4 programmes sur 4 tâches, IQM et intervalle, probabilité d'amélioration par paire. &
160 essais & principale \\
E2, Bonferroni & Le verdict de E2 résiste-t-il à 24 comparaisons ? &
Les 24 intervalles de E2 au niveau simultané de 95~\%. & ceux de E2 &
complémentaire \\
E3 & Deux programmes du même algorithme donnent-ils le même résultat ? &
PPO PyTorch contre PPO JAX, 4 tâches, mêmes graines aléatoires. & 80 de ceux de E2 & principale \\
E4 & Que devient la moyenne quand on garde les 3 meilleurs essais sur 10 ? &
Moyenne et écart type des 10 puis des 3 meilleurs, inversions parmi 24 comparaisons. & ceux de E2 &
principale \\
E5 & Combien de graines aléatoires retrouvent le bon classement ? &
Pour chaque comparaison, \num{10000} tirages avec remise de $n$ essais, $n$ de 3 à 20. &
ceux de E2, les 30 de E1 & principale \\
E6 & Que change une référence réglée selon les réglages publiés ? &
PPO PyTorch sur \task{Hopper}, valeurs par défaut contre réglages du Zoo, puis diagnostic et script
officiel. & 40 essais de plus que E2, plus 10 officiels & principale, diagnostic complémentaire \\
E7 & Que coûte une physique modifiée ? &
Rapport du score en physique modifiée au score en physique d'entraînement, par essai. & ceux de E2 &
principale, classement descriptif \\
Physique par degrés & La perte de E7 tient-elle à une seule modification ? &
Masses de 0,7 à 1,3 fois, frottements de 0,6 à 1,2 fois, une grandeur à la fois. &
320 réévaluations (avec vingt graines aléatoires) & complémentaire \\
B3 & La mesure choisie change-t-elle le classement ? &
Score final, meilleur score, moyenne de la courbe, inversions entre mesures. & ceux de E2 &
complémentaire \\
Vingt graines aléatoires & Les résultats à dix graines aléatoires tiennent-ils ? &
Graines aléatoires 10 à 19, mêmes règles, E2, E3, E5, E7, B3, physique par degrés. & 150 essais de plus &
complémentaire \\
Budget égal & Le classement de E2 tient-il au budget de PPO ? &
PPO lu à \num{200000}~pas, E2 et B3 recalculés. & ceux de E2 &
complémentaire \\
E4 répliquée & Le gonflement de E4 se refait-il ? & E4 à l'identique sur les graines aléatoires 10 à 19. &
160 essais des graines aléatoires 10 à 19 & complémentaire \\
\midrule
E1, seize couples & Le taux de faux positifs de E1 vaut-il pour les autres couples ? &
Les deux tirages de E1, à cinq contre cinq, sur chacun des seize couples. & les 320 de E2 à vingt graines aléatoires &
exploratoire \\
E2 strict & Combien de paires s'inversent vraiment ? &
Paires tranchées dans un sens sur une tâche, dans l'autre sur une autre, sans nouveau tirage. & ceux de E2 &
exploratoire \\
E3 corrigé & La différence de E3 résiste-t-elle à quatre tâches ? &
Mann-Whitney et Holm sur quatre tâches, intervalles à 98,75~\%, Hodges-Lehmann, puissance. & ceux de E3 &
exploratoire \\
E5 hors échantillon & E5 se retrouve-t-il sur d'autres graines aléatoires ? &
Ordre des graines aléatoires 0 à 9 jugé sur les graines 10 à 19, puis l'inverse. & ceux de E2 à vingt graines aléatoires &
exploratoire \\
E7 testé & Le renversement sur \task{Hopper} est-il établi ? &
Probabilité d'amélioration par paire, dans les deux physiques et à chaque degré. & ceux de E2 &
exploratoire \\
Règle calibrée & Quel est le niveau réel de la règle du vainqueur établi ? &
Groupes disjoints d'un même programme, 10 contre 10 et 5 contre 5. & les 30 de E1, les 30 de TRPO, ceux de E2 &
exploratoire \\
E6, vitesse & Le plateau de la recette du Zoo est-il un robot immobile ? &
Journaux d'entraînement, rejeu avec la vitesse, vidéo, agent publié par le Zoo. & ceux de E6 &
exploratoire \\

\end{longtable}
}

\subsection{E1 : même code, autre graine aléatoire}\label{sec:plan-e1}
E1 entraîne trente fois PPO PyTorch sur \task{HalfCheetah-v5}, un million de pas chacun, graines aléatoires 0
à 29, sans rien changer d'autre. Le script \file{sim/analysis/e1.py} lit les trente scores finaux, puis fait,
avec un générateur de graine aléatoire 2018 :
\begin{itemize}
  \item les statistiques de la distribution : moyenne, écart type, médiane, minimum, maximum, coefficient de
    variation (écart type rapporté à la moyenne), test de Shapiro-Wilk ;
  \item pour $k=3$, 5 et 10, \num{10000} tirages de deux groupes disjoints de $k$ essais, obtenus par une
    permutation des trente puis les $k$ premiers contre les $k$ suivants. Chaque tirage reçoit un test de Welch
    et la différence relative des moyennes $|\bar a-\bar b|/\min(\bar a,\bar b)$. Le script compte la part
    des tirages où $p<0{,}05$ et donne la médiane et le 95\textsuperscript{e} centile de la différence relative ;
  \item la coupe \enquote{triée} : \num{10000} tirages de dix essais sans remise, triés par score, puis un
    test de Welch des cinq meilleurs contre les cinq pires. Cette coupe montre ce que produirait un tri des
    essais avant la comparaison ;
  \item la sélection des trois meilleurs sur dix : sur \num{10000} tirages de dix essais, la moyenne des trois
    meilleurs rapportée à la moyenne des dix.
\end{itemize}
Puisque les deux groupes sortent toujours du même programme, toute différence déclarée est un faux positif.
Avec des groupes tirés au hasard et un test bien calibré, la part attendue est de 5~\%.

\paragraph{Les seize couples}
E1 n'étudie qu'un couple, PPO PyTorch sur \task{HalfCheetah}. Une analyse exploratoire refait ses deux tirages à cinq contre cinq sur chacun des seize couples de programme et de tâche de E2,
graines aléatoires 0 à 19 (\file{sim/analysis/e1\_grid.py}) : \num{10000} tirages au hasard, \num{10000} tirages
triés, la part des tirages où le test de Welch donne $p<0{,}05$. Pour juger de la précision de ces parts, le
script rééchantillonne aussi \num{1000} fois les essais de chaque couple avec remise. Les doublons que crée ce
rééchantillonnage écrasent la variance à l'intérieur des groupes et gonflent le taux trié : ces intervalles
décrivent la fragilité de l'estimation, sans garantie de couverture.

\paragraph{La seconde statistique}
La figure~5 de \textcite{henderson2018matters} compare des courbes d'apprentissage entières
(\cref{sec:t-moyen}), là où E1 compare des scores finaux.
\file{sim/analysis/e1\_curve.py} refait leur statistique sur les trente courbes : \num{10000} tirages de dix
essais, coupés en deux groupes de cinq au hasard ou triés par score final, puis un test de Welch à chacun des
vingt points d'évaluation. Le script en tire la part des tirages déclarés différents par le t moyen, par le p moyen
et par le test au dernier point, le 95\textsuperscript{e} centile de $|\bar t|$ et la part des tirages où
$|\bar t|\ge 9{,}09$. Cette statistique complète l'analyse principale de \file{e1.py} sans la remplacer.

\subsection{B1 : l'énigme du t de 9}\label{sec:plan-b1}
\textcite{henderson2018matters} publient, pour TRPO sur \task{HalfCheetah}, $t=-9{,}09$ et $p=0{,}0016$ entre
deux groupes de cinq essais du même code. Pineau précise sur scène que ces groupes ne sont pas triés (faille F11, \cref{sec:failles}). B1 cherche d'où
vient ce chiffre, en deux temps.

\paragraph{Premier temps : leur algorithme}
Le premier temps lance trente essais de TRPO sur
\task{HalfCheetah-v5}, graines aléatoires 0 à 29, un million de pas, avec les valeurs par défaut de
sb3-contrib~2.9.2 et la même évaluation que E1. \file{sim/analysis/b1\_trpo.py} applique les trois lectures de
la \cref{sec:t-moyen} aux courbes de TRPO et de PPO, avec les mêmes tirages que \file{e1\_curve.py}. La règle
de lecture est la suivante. Si le 95\textsuperscript{e} centile de $|\bar t|$ au hasard approche 9 avec
TRPO, la dispersion de l'algorithme explique leur chiffre. Si seule la lecture empilée l'approche, leur calcul
l'explique. Si aucune lecture ne l'approche, le rapport conclut qu'aucune des deux pistes n'explique leur
chiffre.

\paragraph{Second temps : leur script}
Le second temps part de deux faits. Le couple publié $(t,p)=(-9{,}09\,;\,0{,}0016)$ impose environ 3,4~degrés de
liberté en test bilatéral. Il exclut la lecture des points empilés, à environ 198~degrés de liberté, qui
donnerait un $p$ de l'ordre de $10^{-16}$. Il exclut aussi le t moyen converti en $p$, puisqu'un test de Welch
sur cinq contre cinq a au moins 4~degrés de liberté. Seule la moyenne des $p$ reste compatible. Le second fait
est le script de 2017 (\cref{sec:script-henderson}). \file{sim/analysis/b1\_script.py} applique ce calcul aux
trente courbes de TRPO et aux trente de PPO, avec les mêmes tirages (générateur de graine aléatoire 2018, \num{10000}
tirages de dix essais coupés en cinq contre cinq, au hasard puis triés), pour $n=10$ et pour $n=5$. Il calcule
aussi, par programme, le rapport de l'écart type moyen entre épisodes au dernier point à l'écart type des
scores finaux entre graines aléatoires.

\paragraph{Ce qui nous sépare des essais de 2017}
Nos épisodes d'évaluation suivent une politique déterministe, alors que les fichiers de progression de 2017
enregistrent les épisodes d'entraînement, joués par la politique stochastique. Nous avons vingt points
d'évaluation, quand leurs courbes comptent une ligne par itération. Nous ne pouvons pas prouver que ce script a produit leur
figure~5. Son nom (\enquote{de la même distribution}) et sa sortie (\enquote{t=\dots, p=\dots}) le rendent
probable.

\subsection{E2 : le classement change-t-il de tâche en tâche ?}\label{sec:plan-e2}
E2 entraîne les quatre programmes sur les quatre tâches, dix graines aléatoires chacun.
\file{sim/analysis/e2.py} lit le score final de chaque essai, puis, pour chaque tâche :
\begin{itemize}
  \item l'IQM de chaque programme avec son intervalle bootstrap à 95~\% (\cref{eq:iqm,eq:bootstrap}) ;
  \item le classement par IQM ;
  \item pour chacune des six paires, la probabilité d'amélioration du mieux classé sur l'autre, avec son
    intervalle (\cref{eq:poi}), avec le drapeau \enquote{tranchée} si l'intervalle exclut 0,5 ;
  \item un profil de performance : pour 25 seuils $\tau$ de 0 à 1,2, la part des essais dont le score, divisé
    par la meilleure IQM de la tâche, dépasse $\tau$.
\end{itemize}
Le programme \emph{en tête} d'une tâche est celui de plus grande IQM. Le classement \emph{change} si le
programme en tête n'est pas le même sur toutes les tâches. Le changement est \emph{établi} quand le programme
en tête d'une tâche est battu, de façon tranchée, par un autre programme sur une autre tâche. Le script compte
ces changements établis : chaque triplet (tâche où $X$ mène, tâche où $X$ perd, programme qui bat $X$) en
compte un.

\paragraph{Avec la correction de Bonferroni}
\file{sim/analysis/e2\_bonferroni.py} refait les mêmes appels, dans le même ordre, après la même graine aléatoire : les
rééchantillons bootstrap sont identiques à ceux de \file{e2.py}, seul le niveau des centiles change, à
99,79~\%. Il compte les changements établis avec la même définition.

\paragraph{Le décompte strict}
La définition du protocole compte un triplet sans demander si le programme en tête bat de façon tranchée,
sur la tâche où il mène, celui qui le bat ailleurs. \enquote{SAC mène sur \task{HalfCheetah} et perd contre CrossQ sur
\task{Walker2d}} y compte pour un changement établi, même quand l'ordre de SAC et de CrossQ sur
\task{HalfCheetah} n'est pas tranché. Un décompte strict, par paires, la complète
(\file{sim/analysis/e2\_strict.py}) : une paire $(A,B)$ \emph{s'inverse de façon établie} quand $A$ bat $B$ de
façon tranchée sur une tâche et $B$ bat $A$ de façon tranchée sur une autre. Le script relit les intervalles
déjà publiés, à 95~\% et au niveau de Bonferroni, sans nouveau rééchantillonnage, pour les quatre analyses (dix
et vingt graines aléatoires, budget propre et budget égal). Il recopie le décompte du protocole à côté. Le
décompte strict au budget égal demande les intervalles de Bonferroni au budget égal, que \file{e2\_bonferroni.py}
calcule aussi.

\subsection{E3 : deux programmes du même algorithme}\label{sec:plan-e3}
\file{sim/analysis/e3.py} oppose PPO PyTorch à PPO JAX sur chaque tâche, avec les mêmes graines aléatoires et
les scores à un million de pas. Il donne les deux IQM avec leurs intervalles, leur rapport, la
probabilité qu'un essai de PPO JAX batte un essai de PPO PyTorch avec son intervalle, le chevauchement des
intervalles, puis le vainqueur selon la règle de la \cref{sec:regle}.

\paragraph{Quatre tâches, une correction}
E3 teste quatre tâches sans correction. L'analyse à vingt graines aléatoires pose en outre une seconde fois la
même question. L'analyse corrigée traite la première multiplicité et nomme la seconde sans la corriger :
l'analyse à dix graines aléatoires et celle à vingt sont corrigées chacune sur ses quatre tâches, sans correction
entre elles. Le script \file{sim/analysis/e3\_holm.py} calcule pour chaque tâche
le test de Mann-Whitney bilatéral de PPO JAX contre PPO PyTorch, la correction de Holm sur les quatre tâches,
l'intervalle de la probabilité d'amélioration au niveau de 98,75~\%, la différence de Hodges-Lehmann en points de
score (\cref{eq:hl}). Le script estime aussi la puissance du test pour $n$ de 10 à 60 essais par programme
et le plus petit effet détectable (\cref{eq:mde}).

\subsection{E4 : les trois meilleurs sur dix}\label{sec:plan-e4}
\file{sim/analysis/e4.py} calcule, pour chaque programme et chaque tâche, la moyenne et l'écart type des dix
essais, puis des trois meilleurs, enfin le gonflement $\bar x_{\text{3 meilleurs}}/\bar x_{10}-1$. Une comparaison
est une paire de programmes sur une tâche : il y en a 24. Elle \emph{s'inverse} si le signe de la différence
des moyennes change quand chaque programme ne garde que ses trois meilleurs essais. Le script résume le tout
par le gonflement médian, son étendue, la médiane du rapport des écarts types et le nombre d'inversions.

\subsection{E5 : combien de graines aléatoires ?}\label{sec:plan-e5}
\file{sim/analysis/e5.py} part des 24 comparaisons. Pour chacune, le bon ordre est le signe de la différence
des moyennes de tous les essais disponibles. La distribution observée des essais tient lieu de la vraie distribution des scores de l'algorithme. Pour
$n=3$, 5, 10 et 20, on tire \num{10000} fois $n$ essais avec remise pour chacun des deux programmes, puis l'on
compte la part des tirages qui retrouvent le bon ordre :
\begin{equation}\label{eq:e5}
  \hat P_n=\frac{1}{N}\sum_{r=1}^{N}\mathbf{1}\Bigl[\operatorname{sign}\bigl(\bar x^{*(r)}_n-\bar y^{*(r)}_n\bigr)
  =\operatorname{sign}(\bar x-\bar y)\Bigr],\qquad N=\num{10000}.
\end{equation}
Pour PPO PyTorch sur \task{HalfCheetah}, l'analyse à dix graines aléatoires tire parmi les trente essais de E1. Le script
publie la médiane de $\hat P_n$ sur les 24 comparaisons et, parmi elles, sur les comparaisons \emph{proches},
dont la différence relative des moyennes $|\bar x-\bar y|/\max(|\bar x|,|\bar y|)$ reste sous 20~\%. Nous avons corrigé
un défaut du script avant le calcul à vingt graines aléatoires : il prenait pour bon ordre la moyenne des dix premiers essais,
même avec vingt. Il prend désormais la moyenne de toutes les graines choisies. L'analyse à dix graines aléatoires n'a pas
changé d'un octet.

\paragraph{Ce que $\hat P_n$ mesure}
Le bon ordre de l'\cref{eq:e5} est le signe des moyennes de l'échantillon même où les tirages se font, avec
remise. $\hat P_n$ estime donc la stabilité du signe observé quand on refait l'expérience avec $n$ essais tirés
de cet échantillon. Il ne mesure pas la probabilité de retrouver l'ordre vrai des deux algorithmes. Tirer
$n=20$ essais parmi dix observés, dans l'analyse à dix graines aléatoires, est une extrapolation. Une validation hors échantillon
(\file{sim/analysis/e5\_holdout.py}) retire cette circularité. Le bon ordre y vient
d'autres essais, les graines aléatoires 10 à 19. Le script compte les comparaisons dont le signe observé sur les
graines aléatoires 0 à 9 se garde sur les graines 10 à 19. Il tire ensuite $n=3$, 5 et 10 essais sans remise
parmi les graines aléatoires 0 à 9, \num{10000} fois, compte les tirages qui retrouvent l'ordre des graines
aléatoires 10 à 19, puis fait le chemin inverse.

\subsection{E6 : la référence réglée}\label{sec:plan-e6}
E6 met à l'épreuve la faille F7, la référence mal réglée : l'avance qu'une \enquote{nouvelle méthode}
s'attribue quand elle est seule à être réglée. Son premier volet oppose, sur \task{Hopper-v5}, dix essais de PPO
PyTorch avec ses valeurs par défaut (ceux de E2) à dix essais avec les réglages du Zoo (\file{e6.py}).

\paragraph{Le diagnostic}
Le premier résultat a conduit à un diagnostic. \file{e6\_diag.py} compare
cinq groupes de dix essais : valeurs par défaut sur v5, Zoo avec l'erreur de facteur d'actualisation sur v5,
Zoo corrigé sur v5, valeurs par défaut sur v4, Zoo corrigé sur v4. Trois comparaisons sont corrigées par Holm :
l'erreur (Zoo corrigé contre Zoo avec erreur, sur v5), le test loyal sur la version du Zoo (Zoo corrigé contre
valeurs par défaut, sur v4), le même test sur v5. Chaque comparaison rapporte le $p$ de Welch, celui de
Mann-Whitney et la probabilité d'amélioration avec son intervalle.

\paragraph{Le pipeline officiel}
Le pipeline officiel demande si le plafond observé au premier
résultat, un score qui cesse de monter dès les premiers blocs d'entraînement, vient de la recette publiée ou de
notre façon de l'appliquer. Le script \file{sim/zoo\_official.sh} lance dix fois le
script officiel du recueil, \texttt{python -m rl\_zoo3.train -{}-algo ppo -{}-env Hopper-v4 -{}-seed i}, pour
$i$ de 0 à 9, dans l'environnement séparé \file{sim/env-zoo}, sur un fil, sans évaluation pendant
l'entraînement (\texttt{-{}-eval-freq -1}). \file{e6\_official.py} évalue ensuite chaque politique officielle
avec notre protocole : dix épisodes déterministes, graines aléatoires d'évaluation \num{10000} à \num{10009},
normalisation des observations chargée puis figée, sans normalisation de la récompense. Deux comparaisons,
corrigées par Holm, opposent les politiques officielles à notre Zoo corrigé sur v4 puis à
nos valeurs par défaut sur v4. La règle de lecture est la suivante. Si l'officiel dépasse nettement \num{1000},
notre application était fautive. S'il plafonne aussi, la recette publiée fait moins bien que les valeurs par
défaut dans cette version des bibliothèques. Le rapport le constate alors, sans parler d'une
référence que nous aurions mal réglée.

\paragraph{Un robot qui tient debout sans avancer}
Sur \task{Hopper}, la récompense d'un pas vaut la vitesse vers l'avant, plus 1 tant que le robot tient debout,
moins $10^{-3}\,\lVert a\rVert^2$ pour l'action $a$. Un épisode de \num{1000}~pas qui rapporte environ
\num{1000} est donc un robot debout et immobile. Trois analyses lisent la vitesse au lieu du score :
\begin{itemize}
  \item la première lit les journaux d'entraînement des dix agents officiels (\file{sim/analysis/e6\_speed.py}) :
    sur les 200 derniers épisodes, la part d'épisodes de \num{1000}~pas et leur retour par pas, dont on déduit
    une borne basse de la vitesse moyenne, le retour par pas moins 1 ;
  \item la deuxième rejoue les politiques officielles, nos politiques réglées comme le Zoo et nos politiques
    par défaut avec notre protocole, en relevant la vitesse et la position à chaque pas, puis filme un épisode
    (\file{sim/analysis/e6\_enjoy.py}). Contrôle : les retours rejoués doivent redonner au chiffre près les
    scores déjà publiés ;
  \item la troisième télécharge l'agent publié par le Zoo pour \task{Hopper-v3}, inspecte chaque
    objet sérialisé avant de le charger, puis l'évalue avec notre protocole sur \task{Hopper-v4}
    (\file{sim/analysis/e6\_hub.py}).
\end{itemize}
Un réentraînement de la recette du Zoo avec $\gamma=0{,}99$ pour l'algorithme et la normalisation, rien
d'autre, dix graines aléatoires sur \task{Hopper-v4}, est analysé par \file{sim/analysis/e6\_gamma.py},
écrit avant son entraînement : moyenne interquartile et son intervalle, probabilité d'amélioration contre
la recette à 0,999 et contre les valeurs par défaut, correction de Holm sur les deux comparaisons, part des
épisodes debout et immobiles.

\subsection{E7 : la physique modifiée}\label{sec:plan-e7}
\file{sim/analysis/e7.py} calcule, pour chaque essai $s$, le rapport des deux scores finaux
\begin{equation}\label{eq:e7}
  \rho_s=\frac{\texttt{perturbed\_mean}_s}{\texttt{eval\_mean}_s},
\end{equation}
puis, pour chaque programme et chaque tâche, l'IQM des $\rho_s$
avec son intervalle bootstrap. Le verdict est une \emph{perte} si l'intervalle reste sous 1, un \emph{gain}
s'il reste au-dessus, \enquote{aucun changement détectable} sinon. L'IQM des rapports diffère du rapport
des IQM : la part de score gardée ne se lit pas sur deux classements de scores.

\paragraph{Le classement descriptif}
\file{e7\_rank.py} classe les quatre programmes sur chaque tâche par IQM du score, dans chacune des
deux physiques. Il dit si le programme en tête change. Ce complément est descriptif, sans test.

\paragraph{Le classement en physique modifiée, testé}
\file{sim/analysis/e7\_test.py} donne un test à ce classement. Pour chaque tâche et chaque paire, le script calcule la probabilité d'amélioration avec son intervalle et la valeur $p$ de
Mann-Whitney, en physique d'entraînement comme en physique modifiée. Une paire \emph{se renverse de façon
établie} quand elle est tranchée dans un sens dans la première et dans l'autre sens dans la seconde. La question
principale est PPO JAX contre CrossQ et contre SAC sur \task{Hopper}. Le reste est descriptif.
Le script refait le même calcul à chaque degré de la physique par degrés.

\paragraph{La physique par degrés}
\file{sim/analysis/dose.py} lit les 160 fichiers de réévaluation (320 avec vingt graines aléatoires). Pour chaque
programme, tâche, grandeur et facteur, il calcule l'IQM du rapport au score à 1,0 fois, avec son intervalle
bootstrap, puis le même verdict que E7. Pour chaque tâche et chaque physique, il classe les quatre programmes par
IQM du score, sans test. Il vérifie enfin que le score à 1,0 fois retrouve le score final de \file{run.py} et note
la plus grande différence relative entre les deux. Le script a été écrit alors que 30 des 160 réévaluations
existaient déjà, dont une avait été affichée. Un commit le disait écrit avant tout résultat : nous avons
corrigé cette déclaration. Le script n'a aucun paramètre libre, si bien que ce qui avait été vu ne pouvait
rien changer à son calcul.

\subsection{B3 : la mesure change-t-elle le classement ?}\label{sec:plan-b3}
\file{sim/analysis/b3.py} calcule trois mesures par essai sur les courbes enregistrées, sans aucun
entraînement :
\begin{itemize}
  \item le score final, au dernier point ;
  \item le meilleur score atteint le long de la courbe ;
  \item la moyenne de la courbe, l'aire sous la courbe divisée par sa longueur.
\end{itemize}
Pour chaque tâche et chaque mesure, il calcule l'IQM de chaque programme et le classement. Une paire de
programmes \emph{s'inverse} entre deux mesures si l'ordre de leurs IQM change. L'inversion est \emph{établie}
si, sous chacune des deux mesures, l'intervalle de la probabilité d'amélioration exclut 0,5, dans des sens
opposés. Les courbes de PPO ont vingt points sur un million de pas, celles de SAC et de CrossQ quatre points
sur \num{200000} : la moyenne de la courbe ne couvre pas le même budget.

\paragraph{Les courbes elles-mêmes}
Trois lectures exploratoires des courbes complètent B3
(\file{sim/analysis/curves.py}) :
\begin{itemize}
  \item le classement des quatre programmes à \num{50000}, \num{100000}, \num{150000} et \num{200000}~pas,
    comparé au classement final par le $\tau$ de Kendall, avec la corrélation de Spearman entre le score d'une
    graine aléatoire en cours de route et son score final ;
  \item la part des essais dont le score final reste sous leur meilleur point, avec le gonflement de l'IQM
    quand on rapporte ce meilleur point ;
  \item l'autocorrélation des courbes et leur nombre effectif de points indépendants (\cref{sec:autres-tests}).
\end{itemize}

\subsection{La consolidation : vingt graines aléatoires et budget égal}\label{sec:plan-consolidation}
La consolidation ajoute trois volets avec le même code et les mêmes réglages :
\begin{itemize}
  \item les graines aléatoires 10 à 19 pour les quatre programmes sur les quatre tâches, soit 150 nouveaux
    entraînements, puisque PPO PyTorch sur \task{HalfCheetah} avait déjà les siens ;
  \item la physique par degrés de la \cref{sec:plan-e7} ;
  \item la lecture à budget égal de la \cref{sec:budgets}.
\end{itemize}
Le module \file{sim/analysis/runs.py} lit trois variables d'environnement, toutes inactives par défaut :
\texttt{SEEDS} (par exemple \texttt{0:20}), \texttt{BUDGET=equal} et \texttt{SUFFIX}, ajouté au nom du fichier
de sortie. Les calculs sont tous lancés par \texttt{make analysis} :
\begin{itemize}
  \item sur les vingt graines aléatoires (suffixe \file{\_s20}) : E2, E3, E5, E7, B3, la physique par degrés, le
    classement descriptif de E7 et la correction de Bonferroni de E2 ;
  \item E4 sur les seules graines aléatoires 10 à 19 (suffixe \file{\_r10}), comme réplication indépendante ;
  \item E2 et B3 à budget égal, sur dix graines aléatoires (\file{\_equal}) et sur vingt (\file{\_equal\_s20}).
\end{itemize}
Les deux derniers calculs de la première liste appliquent aux vingt graines aléatoires des règles déjà
écrites. Les résultats à dix graines aléatoires restent les résultats principaux.

\section{La file d'attente et le calcul}\label{sec:calcul}

\paragraph{La machine}
Tous les entraînements ont tourné sur processeur, sur le Mac de la \cref{sec:outils}, partagé avec d'autres
projets. Un entraînement ou une réévaluation occupe un cœur (un seul fil de calcul). Le
\cref{tab:calcul} donne, par programme et par tâche, le nombre d'essais, les pas, la durée médiane d'un essai et
la somme des durées, en heures.

\gentable{calcul}

\paragraph{Les lanceurs, dans l'ordre où ils ont servi}
Plusieurs scripts ont lancé les essais. Chacun saute un essai dont la politique finale existe déjà.
\begin{itemize}
  \item \file{sim/drive.sh} lit une liste de tâches (\file{jobs.txt} par défaut) et en lance $P$ à la fois par
    \texttt{xargs -P}. Un essai qui échoue est relancé une fois, dans un journal \file{.retry.log}.
    \file{drive2.sh} en est une copie pour \file{jobs2.txt}, les essais de E6 corrigés.
  \item \file{sim/relay.sh} change de file sans tuer d'essai : il attend la fin des essais en cours, puis
    relance la file principale à un entraînement JAX (environ deux cœurs) et la file TRPO de B1 à deux
    entraînements, pour rester à quatre cœurs.
  \item \file{sim/zoo\_official.sh} attend la fin de la file TRPO, puis lance les dix entraînements officiels de
    E6 sur deux voies de cinq.
  \item \file{sim/queue.sh} lit \file{jobs\_consolidation.txt}, la liste de la consolidation : 160 lignes
    \texttt{train} et 320 lignes \texttt{dose}. Il réserve chaque tâche par un \texttt{mkdir} atomique, ce qui
    permet à deux copies de tourner sans se gêner. Il tient le registre \file{logs/attempts.jsonl}. Un arrêt par
    signal, code de sortie de 128 ou plus, n'est pas relancé.
  \item \file{sim/worker.sh} est \file{queue.sh} avec un plafond : il lit \file{workers.max} avant chaque tâche
    et s'en va si trop de copies tournent.
  \item \file{sim/scale.sh} décide ce plafond toutes les heures. Il somme le calcul pris par les autres
    programmes de la machine (\texttt{\%cpu} de \texttt{ps}), en déduit les cœurs libres moins un de marge. Il
    accorde au projet au plus \texttt{MAX\_CORES} cœurs, l'essai GRPO sur processeur en prenant deux. Il lance
    de nouveaux \file{worker.sh} si le plafond monte. Il n'en lance aucun sans le surveillant. Ses décisions
    sont dans \file{logs/scale.log}.
  \item \file{sim/watchdog.sh}, le surveillant, mesure toutes les 20~minutes l'espace disque libre et la
    croissance de la mémoire d'échange (\emph{swap}) depuis son lancement. Il arrête tous les lanceurs et tous
    les essais si le disque passe sous \texttt{MIN\_FREE\_GB} gigaoctets (10 par défaut) ou si la mémoire
    d'échange a grossi de plus de 8~Go. Il inscrit chaque tâche arrêtée au registre et libère sa réservation.
\end{itemize}

\paragraph{Le nombre de cœurs}
Le projet a pris de deux à huit cœurs selon la charge de la machine : quatre réglés à la main pour la
grille principale, puis deux copies de \file{queue.sh} pour la consolidation, puis un plafond décidé
chaque heure par \file{scale.sh}, de huit cœurs au plus, enfin deux cœurs pour tout le projet, essai GRPO
sur processeur compris. Chaque essai tourne sur un seul fil, dans son
propre processus : le nombre de tâches parallèles n'entre dans aucun de ses calculs. Les relances de PPO
retrouvent leurs points d'évaluation à l'identique. Pour SAC et CrossQ, trois entraînements refaits à
l'identique aussi.

\paragraph{Les arrêts et les relances}
Le registre reconstruit \file{logs/attempts\_reconstructed.csv} compte 292~lancements pour les
250~entraînements antérieurs à la consolidation : 42~tentatives arrêtées, aucune plantée. Le protocole
demande de noter chaque incident. Seules 2 de ces 42 tentatives l'avaient été : nous avons corrigé ce
défaut par ce registre, puis par celui que tient \file{queue.sh}. Les arrêts ont trois causes : un premier arrêt général, probablement celui de l'incident
mémoire du pilote GRPO, un arrêt demandé pour une machine surchargée, puis quatre arrêts pour
rester à quatre cœurs. Chaque arrêt a frappé tous les essais en cours à la fois. Aucun essai n'a été relancé ni
écarté après que son score final était connu, ni pour une raison liée à son score. Pour PPO PyTorch et PPO
JAX, les relances retrouvent tels quels les points d'évaluation des tentatives arrêtées. Pour SAC et
CrossQ, les courbes arrêtées ont été écrasées : cette vérification est impossible. Les arrêts suivants
du surveillant sont inscrits au registre \file{logs/attempts.jsonl}.

\paragraph{La chaîne jusqu'aux documents}
Les essais écrivent leurs fichiers CSV et JSON. Les scripts de \file{sim/analysis/} en tirent les fichiers
JSON de résultats. \texttt{make verify} recalcule chaque analyse et compare les fichiers à l'octet près, puis
\file{docs/check\_numbers.py} et \file{article/check\_article.py} vérifient chaque nombre des documents contre
ces fichiers. L'annexe~\ref{chap:08-logiciel} décrit cette chaîne et la dessine (\cref{fig:architecture,fig:architecture-2}).
